%----------------------------------------------------------
%	PACKAGES AND THEMES
%----------------------------------------------------------
\documentclass[aspectratio=169,xcolor=dvipsnames,presentation]{beamer}
%\documentclass[aspectratio=169,xcolor=dvipsnames,handout]{beamer}

\usetheme{Darmstadt}
\usecolortheme{seahorse}
\usepackage[hangul]{kotex}
\usepackage{hyperref}
\usepackage{graphicx, array, adjustbox }
\usepackage{xcolor, transparent}
\usepackage{booktabs, multicol, multirow}
\setbeamercovered{transparent}

%----------------------------------------------------------
%	TITLE PAGE
%----------------------------------------------------------

\title[]{고용관계에 대하여} % The short title appears at the bottom of every slide, the full title is only on the title page
\subtitle{노사관계의 이론과 실제}
\author[]{오성재}
\institute[CNU] % Your institution as it will appear on the bottom of every slide, may be shorthand to save space
{%
    충남대학교 경제학과\\
}
\date{2024년 9월 4일} % Date, can be changed to a custom date

\begin{document}

\begin{frame}
    \titlepage%
\end{frame}

\begin{frame}{목차}
    \tableofcontents[hideallsubsections]
\end{frame}

\section{고용관계의 배경}

\subsection{정보화·세계화와 고용관계}

\begin{frame}{표}
    \scriptsize
    \begin{table}
        \centering
        \resizebox{.6\textwidth}{!}{%
            \begin{tabular}{|c|c|}
\toprule
\textbf{항목 1} & \textbf{항목 2} \\ \midrule
\textbf{내용 1} & \textbf{내용 2} \\ \hline
\textbf{내용 3} & \textbf{내용 4} \\ \bottomrule
\end{tabular}
        }
        \\
        \raggedright % 왼쪽 정렬
        \hspace{1em} % 좌측 여백을 추가
        \tiny{자료:} % 자료 출처
        \caption{표제목}
    \end{table}
\end{frame}

\begin{frame}{그림2단}
\begin{columns}
    \begin{column}{.5\textwidth}
        \begin{figure}
            \centering
            \includegraphics[width=.7\textwidth]{pic/smileman.png}
            \\
            \raggedright % 왼쪽 정렬
            \hspace{1em} % 좌측 여백을 추가
            \tiny{자료: } % 자료 출처
            \caption{좌측그림}
        \end{figure}
    \end{column}
    \begin{column}{.5\textwidth}
        \begin{figure}
            \centering
            \only<1|handout:0>{\transparent{0.3}\includegraphics[width=.7\textwidth]{pic/smileman.png}} % 첫 번째 슬라이드: 투명하게 표시
            \only<2>{\includegraphics[width=.7\textwidth]{pic/smileman.png}} % 두 번째 슬라이드: 선명하게 표시
            \\
            \raggedright % 왼쪽 정렬
            \hspace{1em} % 좌측 여백을 추가
            \tiny{자료:} % 자료 출처
            \caption{우측그림}
        \end{figure}
    \end{column}
\end{columns}
\end{frame}

\begin{frame}{큰 그림}
\centering
\begin{figure}
    \includegraphics[width=.3\textwidth]{pic/smileman.png}
    \\
    \raggedright % 왼쪽 정렬
    \hspace{1em} % 좌측 여백을 추가
    \tiny{자료: } % 자료 출처
    \caption{제목}
\end{figure}
\end{frame}

\begin{frame}{그림글 2단}
\begin{columns}
    \begin{column}{.5\textwidth}
        \begin{figure}
            \centering
            \includegraphics[width=.7\textwidth]{pic/smileman.png}
            \\
            \raggedright
            \hspace{1em}
            \tiny{자료:}
            \caption{그림제목}
        \end{figure}
    \end{column}

    \begin{column}{.5\textwidth}
        \uncover<1->{%
        \begin{itemize}
            \item \uncover<2->{나랏말싸미}
            \item \uncover<3->{나랏말싸미}
        \end{itemize}
        }
    \end{column}
\end{columns}
\end{frame}

\begin{frame}{경고블록}
    \begin{itemize}[<+->]
        \item 나랏말싸미
        \begin{itemize}[<+->]
            \item 나랏말싸미
        \end{itemize}
    \end{itemize}
    \begin{alertblock}{경고블럭}
        \begin{itemize}[<+->]
            \item 나랏말싸미
        \end{itemize}
    \end{alertblock}
\end{frame}

\begin{frame}{예시블록}
    \begin{itemize}[<+->]
        \item 나랏말싸미
        \begin{itemize}[<+->]
            \item 나랏말싸미
        \end{itemize}
    \end{itemize}
    \begin{exampleblock}{블록제목}
        \begin{itemize}[<+->]
            \item 나랏말싸미
        \end{itemize}
    \end{exampleblock}
\end{frame}

\begin{frame}{블록}
    \begin{itemize}[<+->]
        \item 나랏말싸미
        \begin{itemize}[<+->]
            \item 나랏말싸미
        \end{itemize}
    \end{itemize}
    \begin{block}{블록제목}
        \begin{itemize}[<+->]
            \item 나랏말싸미
        \end{itemize}
    \end{block}
\end{frame}

\begin{frame}{개조식}
    \begin{itemize}[<+->]
        \item 나랏말싸미
        \begin{itemize}[<+->]
            \item 나랏말싸미
        \end{itemize}
    \end{itemize}
\end{frame}

%------------------------------------------------
\end{document}
%------------------------------------------------
